\section{Information Heist}
\label{sec:informationHeist}

The party may decide to break into \locationIndicium{} and steal the information they need. In that case, run the scene as a Mission Endeavor.

During the Mission Endeavor, the PCs may pursue multiple objectives. Each additional objective adds more successes required to complete the Endeavor, but the number of acceptable failures does not change. 

The PCs may decide which objective are they pursuing with each of their tests. Once they reach the number of successes required to complete the objective, they earn the respective reward and may continue with other objectives.




\subsection{Mission Goals}

During the Mission, the PCs may work towards the goals listed below. They may also define their own objectives - in that case, use your own judgment how many successes should they require.

\begin{itemize}
	\item Finding the location of Singer informant - 1 success per PC
	\item Leaving \locationIndicium{} without being captured - 0.5 of success per PC
	\item Learning the story of the orphanage - 0.5 of success per PC
	\item Fining clues about the Cult's hideout - 0.5 of success per PC
	\item Stealing top-secret state information - 1 success per PC
\end{itemize} 

\reasonBox{
	This book counts successes per PC to better balance for parties with different size than 4. This way, each PC's contribution is important when the party is bigger and the players are not stuck in endless dice rolling in smaller parties. If you don't like that approach, multiply the values above times 4, as this is the standard party size the book assumes. 
}

\reasonBox{
	If the party has an odd number of PCs, the partial goals will not be a whole numbers. In that case round them up after summing all of them. Example: if a party of 3 PCs aims to learn about the singer, orphanage and leave safely, they will need to gather 6 successes in total (3 + 1.5 + 1.5 = 6). 
}




\subsection{Available Time}

The party has limited time to complete the Mission. In total, they may do 3 tests per PC. Raise stakes with each test. On \iconComp{} you may reduce that pool by 1 available test, reducing the total time of the Mission.

\ptrBox{
	Remaining in control while in stressful situation may be recognized by an ashspren.
}




\subsection{Mission End}

Once the party runs out of time, start \fullref{sec:wanted} event unless it is already in place. Then advance the event by one \iconComp{}. During the scene you may spend multiple \iconComp{} to progress the event.

The mission ends when the PCs are captured or after they complete "Leave \locationIndicium{}" goal. At this point, they cannot proceed with more tests to achieve more goals.




\subsection{Example Tests}

During the Mission, the PCs may attempt one of the below suggestion or propose their own ideas. The DC of the tests should fall in the range of 13-16 (12 + 1d4) on tier 2, or (8 + 1d4 + 2 $\times$ tier) on other tiers.

\scriptedOption{Quiet sneaking}{
	A PC may attempt a \skillCheck{15}{\skillStealth} to sneak past \enemyGuardB{}.
}

\scriptedOption{Lie to a scribe}{
	When spotted by a \enemyScribeB{}, a PC may attempt \skillCheck{16}{\skillDeception{}} to convince them they are working here.
}

\scriptedOption{Lockpicking}{
	A PC may attempt a \skillCheck{14}{\skillThievery} to access a room behind locked door.
}

\scriptedOption{Archive system}{
	A PC may attempt \skillCheck{16}{\skillDeduction} to determine the system in which the scribes are archiving the data.
}

\scriptedOption{Fast reading}{
	To quickly scan large amounts of text, a PC may attempt a \skillCheck{14}{\skillPerception}.
}

\scriptedOption{Useful contraption}{
	A PC may try to craft something helpful (a mirror to peek around the corners, perhaps?) with a \skillCheck{13}{\skillCrafting}.
}


\subsection{Opportunities and Complications}

During the Mission, the following examples might be used when \iconOpp{} or \iconComp{} is rolled:

\scriptedOption{Time's up!}{
	A GM may spend \iconComp{} to reduce the number of available skill test attempts by 1. Someone just spotted something suspicious, hears the party, or takes a wrong turn.
}

\scriptedOption{Determined to succeed}{
	A player may spend \iconOpp{} from a failed test to repeat the test.
}

\scriptedOption{How could you...?}{
	The GM may spend \iconComp{} to reveal a secret of one of the PCs to another PC as they found a report related to it.
}

\scriptedOption{Curious trivia.}{
	Spend \iconOpp{} to reveal a curious secret that the PCs did not know they could search for.
}

\scriptedOption{Sidetracked}{
	Spend \iconOpp{} or \iconComp{} to decrease or increase a number of required successes of uncompleted secondary goals (other than Finding Singer or Leaving Undetected).
}

\scriptedOption{Path to Radiance}{
	\iconOpp{} or \iconComp{} can be spent to increase or reduce the progress of one of personal goal of a PC as they find some useful/demeaning information.
}



\subsection{Secrets to Find}

During the Mission, the party may find multiple secrets in \locationIndicium{}, either by completing their objectives, or by stumbling on one by accident. You can use the secrets from the list below to share with the PCs.

\subsubsection{Singer in \locationYeddaw}

If the PCs complete the main objective, they will find the following report:

\quoteBox{
	TODO
}

It seems that \npcInformant{}, the Singer informant, can be found in the Immigrant Quarter. The party should look for him there - the PCs can now proceed to \fullref{sec:searchForSinger}.

\subsubsection{Tashi's Light Orphanage}

If the PC find the reports of the orphanage, read the following:

\quoteBox{
	TODO
}

\npcArshqqam{} manifested magical powers and accepted an invitation to the emperor palace in Azimir. Without her, surgeons from \locationKharbranth{} took over the orphanage, possibly with hopes to convert it to a hospital.

\subsubsection{State Secrets}

If the PCs uncover this secret, read the following:

\quoteBox{
	TODO
}

The party learns that the prince of \locationTashikk{} worked with \npcNale{}, the Herald of Justice, to execute all power wielders and prevent the return of the Desolations. The recent red storm proves the Herald failed.

This secret does not help the PCs to reach their goal, but perhaps someone will be willing to trade for it.

\subsubsection{Curious Little Secrets}

During the search, the party may learn some small, curious, but potentially irrelevant secrets:

\begin{itemize}
	\item The emperor is a good friend with a power wielder kid, one can expect power-friendly policies from Azir.
	\item A ship carrying a shardbearer was sunk during the red storm while on the voyage to Icewater.
	\item New fashion trends from Liafor, designed for the upcoming year.
	\item The most recent map of Reshi Isles and their predicted movement.
	\item Confirmed: the Assassin in White was slain by the alethi during the recent battle.
	\item New fabrial design from Jah Keved.
	\item Rira plans to sow more discord on Iri territory. \textit{(detailes plans included)}
\end{itemize}

\reasonBox{
	If you are planning more adventures after this one, you can use this opportunity to foreshadow or plot hook the future events.
}